% =============================================================================
% Kapitel 6: Diskussion
% =============================================================================
\chapter{Diskussion}
\label{ch:diskussion}

Dieses Kapitel diskutiert die Ergebnisse der Evaluation (\cref{ch:evaluation}) im Kontext der Forschungsfragen und der verwandten Arbeiten. Es werden die Limitationen der Arbeit aufgezeigt und die Implikationen der Ergebnisse für die Praxis und die Forschung erläutert.

% -----------------------------------------------------------------------------
\section{Interpretation der Ergebnisse}
\label{sec:interpretation}

\subsection{FF1 -- Zuverlässige Vorhersage des individuellen Mobilitätsverhaltens} 

FF1 wird durch die Ergebnisse gestützt, sofern hochwertige Daten vorliegen; \\\texttt{yjmob100k} liefert hierfür den aussagekräftigsten Beleg, da die Quelle reale Bewegungsdaten von 50 Individuen umfasst und somit über Personen hinweg generalisiert. Der Random Forest erreicht dort eine PR-AUC von $0,683$ (\cref{tab:baselines}), während die reine Zufalls-Baseline bei $0,382$ liegt. Ausschlaggebend ist letztlich nicht der Abstand zum Zufall, sondern der Abstand zur stärksten \emph{informierten} Baseline -- also die stärkste Metrik, welche keine komplexe Berechnung oder Algorithmen benötigt -- die Klimatologie, die allein Wochentag und Stunde nutzt, erreicht bereits $0,489$. Der Random Forest übertrifft sie um $0,194$ und damit um $40\%$. Da die Streuung über die fünf Seeds mit $\sigma \approx 0,0004$ gegenüber diesem Abstand um mehr als zwei Größenordnungen kleiner ist, ist der Vorsprung statistisch eindeutig und kein Artefakt der Zufallsinitialisierung. Die PR-AUC dient als Leitmaß, da sie gegenüber der Accuracy robuster ist, insbesondere bei niedrigen Aktivraten. Bei \texttt{emobpy} beispielsweise schlägt der triviale \enquote{immer geparkt}-Klassifikator das Modell sogar ($0,904$ gegenüber $0,697$, \cref{tab:baselines}), ohne eine einzige Fahrt zu erkennen. Der Befund belegt somit, dass das Modell, jenseits der reinen Tages- und Wochenperiodik, Struktur nutzt -- genau das, was FF1 unter einem \enquote{ausreichend komplexen Modell} versteht. Die Gültigkeit bleibt jedoch an die Datenqualität gebunden: Auf \texttt{ved} erreicht dasselbe Verfahren nur $0,089$ und liegt damit praktisch auf dem Niveau seiner Klimatologie-Baseline ($0,083$), sodass FF1 dort nicht aussagekräftig beantwortet werden kann.


\subsection{FF1.1 -- Unterschiede der Machine Learning Verfahren in Prognosegenauigkeit}

Die unterschiedlichen Verfahren zeigen in der Evaluation (\cref{ch:evaluation}), dass die Modelle sich in ihrer Leistung nicht wesentlich unterscheiden. Die Experimente zeigen, dass bestimmte Modellarchitekturen nicht für jeden Datensatz geeignet sind. So erzielt LightGBM auf \texttt{real\_world\_ev} einen F1-Score von $0,366\pm0,0$ und PR-AUC von $0,256\pm0,0$, während Random Forest dort $0,680\pm0,037$ und $0,630\pm0,013$ erreicht. Andersrum zeigt sich auf \texttt{ved}, dass LightGBM mit einem F1-Score von $0,092\pm0,0$ erreicht, während Random Forest nur $0,075\pm0,010$ erreicht; bei den restlichen Werten dominiert jedoch wieder der Random Forest (\cref{tab:cv_metrics}). Die LSTM-Modelle lassen sich auf der Mehrheit der Datensätze nicht abbilden, da diese die Mindestanforderung von 100 Fahrzeugwochen nicht erfüllen. Die Ergebnisse der LSTM-Varianten sind daher nur eingeschränkt interpretierbar, wobei das tabellarische LSTM bei \texttt{yjmob100k} bei F1-Score ($0,677\pm0,001$) und PR-AUC ($0,702\pm0,001$) vorne liegt und das sequentielle LSTM bei \texttt{emobpy} bei F1-Score ($0,324\pm0,003$), ROC ($0,820\pm0,002$) und PR-AUC ($0,289\pm0,006$). Insgesamt erzielen die Machine Learning Verfahren jedoch auf allen Datensätzen bessere Ergebnisse als die Baselines, sodass FF1.1 im Sinne der Forschungsfrage beantwortet werden kann.

\subsection{FF1.1.1 -- Einfluss der Datenquelle auf die Prognosegenauigkeit}

Aufbauend auf FF1.1 zeigt sich, dass die Datenquelle einen entscheidenden Einfluss auf die Prognosegenauigkeit hat. Die Ergebnisse der Evaluation (\cref{ch:evaluation}) zeigen, dass die Wahl der Datenquelle ein größerer Faktor für die Prognosegenauigkeit ist als die verwendeten Machine Learning Verfahren. In \cref{tab:cv_metrics} zeigt sich beispielsweise, dass der Random Forest auf \texttt{yjmob100k} einen F1-Score von $0,669\pm0,001$ erreicht, während er auf \texttt{ved} nur $0,075\pm0,010$ erreicht. Im Vergleich schneidet LightGBM mit dem F1-Score bei \texttt{yjmob100k} mit $0,676\pm0,0$ ab, während es bei \texttt{ved} nur $0,092\pm0,0$ sind. Die Unterschiede zwischen den Datensätzen sind damit größer als die Unterschiede zwischen den Verfahren. Dies deutet darauf hin, dass die Qualität und Repräsentativität der Datenquelle entscheidend für die Prognosegenauigkeit ist. Dadurch lässt sich die Forschungsfrage FF1.1.1 beantworten: Die Datenquelle hat einen signifikanten Einfluss auf die Prognosegenauigkeit, so sehr, dass die Wahl der Datenquelle wichtiger ist als die Wahl des Machine Learning Verfahrens.

\subsection{FF1.2 -- Merkmalseinfluss auf individuelles Mobilitätsverhalten}

In \cref{tab:ablation} zeigt sich, dass gewisse Merkmalsgruppen einen größeren Einfluss auf die Prognosegenauigkeit haben als andere. Die Merkmalsgruppe \texttt{Lag} zeigt den größten Einfluss auf die PR-AUC Metrik, während \texttt{Rolling} teilweise sogar einen negativen Einfluss auf die Prognosegenauigkeit hat. In \cref{tab:ablation_rf} ist besonders bei \texttt{real\_world\_ev} für den F1-Score zu erkennen, dass die Merkmalsgruppe \texttt{Lag} den entscheidenden Faktor darstellt, denn die Prognose bricht ohne diese Merkmalsgruppe auf $0,067$ zusammen, die Prognose ist mit allen Merkmalen bei $0,619$, während der Wert mit \texttt{Lag} als einzige Merkmalsgruppe bei $0,669$ liegt. Der einzige Datensatz, bei dem das Wegnehmen der Merkmalsgruppe \texttt{Lag} keinen signifikanten Einfluss auf die Prognosegenauigkeit hat, ist \texttt{ved}, was jedoch daran liegt, dass die Prognosegenauigkeit dort ohnehin sehr niedrig ist. Die Merkmalsgruppe \texttt{Rolling} sorgt bei drei Datensätzen für eine Verschlechterung der Prognosegenauigkeit, was darauf hindeutet, dass diese Merkmalsgruppe nicht für alle Datensätze geeignet ist. Damit kann die Forschungsfrage beantwortet werden: Die Merkmalsgruppe \texttt{Lag} ist somit die wichtigste Merkmalsgruppe für die Vorhersage des individuellen Mobilitätsverhaltens, während \texttt{Rolling} nur in bestimmten Fällen einen positiven Einfluss hat. 

\subsection{FF2 -- Anforderung an Daten}

Die in \cref{tab:anforderungen} zusammengefassten Anforderungen an die Daten werden durch die Evaluation (\cref{ch:evaluation}) gestützt. Durch die Ergebnisse der Evaluation wird deutlich, dass die Qualität und Repräsentativität der Datenquelle entscheidend für die Prognosegenauigkeit ist. Die Anforderungen an die Daten sind daher gerechtfertigt und notwendig, um eine zuverlässige Vorhersage des individuellen Mobilitätsverhaltens zu ermöglichen. Eigenschaften eines Datensatzes, die notwendig sind, um eine zuverlässige Vorhersage des individuellen Mobilitätsverhaltens zu ermöglichen, sind insbesondere ein ableitbares Label für die Fahraktivität, eine ausreichende Anzahl an Fahrzeugwochen um dem Modell genügend Trainingsdaten zu liefern, sowie eine handhabbare Imbalance der Klassenverteilung, sodass das Modell nicht nur die Mehrheit der Daten lernt. Ein weiteres Kriterium, das durch \texttt{ved} bewusst wurde, ist die zeitliche Verteilung und Zeitspanne von Individuen. Teilweise sind Fahraktivitäten, durch ihre begrenzte zeitliche Verteilung, nur im Trainingsdatensatz enthalten, sodass das Modell diese nicht testen und lernen kann. 

\subsection{FF2.1 -- Aufbereitung und Anonymisierung}

Im Vergleich steht \texttt{yjmob100k}, dessen Aktivitätserfassung durch Smartphone GPS-Daten und durch das Überschreiten von Rastern mit 500Meter Kantenlänge erfolgt. Anders als bei \texttt{yjmob100k} notieren die anderen Datensätze die Fahraktivität durch die Existenz der Einträge bei \texttt{ved} oder ein konkretes Label bei \texttt{real\_world\_ev}. Die Evaluation (\cref{ch:evaluation}) zeigt, dass die Art der Aufbereitung und Anonymisierung der Daten einen Einfluss auf die Prognosegenauigkeit hat. So ist es schwierig bei \texttt{yjmob100k} eine überragende Leistung zu erzielen, da die Fahraktivität nur durch das Überschreiten von Rastern erfasst wird und somit eine gewisse Ungenauigkeit in der Labelerstellung besteht. Folglich lässt sich die Forschungsfrage FF2.1 beantworten: Die Art der Aufbereitung und Anonymisierung der Daten hat einen Einfluss auf die Prognosegenauigkeit, da sie die Qualität der Labels beeinflusst.

\subsection{FF2.2 -- Bewertung der Datenqualität}

Die Qualität der Daten determiniert die Güte der Modelle. Wie man in der \cref{tab:cv_metrics} und \cref{tab:ablation} sehen kann, gibt es einen Datensatz, der sich deutlich von den anderen abhebt: \texttt{ved}; und zwar durch seine besonders schlechten Modellleistungen. Das Problem ist hier spezifisch, dass Daten nur willkürlich auftauchen oder Fahrzeuge teils nur in bestimmten Zeitslots aktiv sind, was zu einem Fehlen in den Test- oder Trainingssets führen kann und somit einige Einträge nicht in beiden Sets repräsentiert sind. Die Ergebnisse zeigen, dass die Qualität der Daten einen entscheidenden Einfluss auf die Prognosegenauigkeit hat. Forschungsfrage 2.2 kann somit wie folgt beantwortet werden: Die Qualität der Daten bildet das Fundament der darauffolgenden Modellierung und ist wegbereitend für ein erfolgreiches Experiment, somit kann ohne eine gute Datenqualität keine zuverlässige Prognose erreicht werden.

% -----------------------------------------------------------------------------
\section{Vergleich mit verwandten Arbeiten}
\label{sec:vergleich}

Ein direkter Vergleich mit bestehenden Arbeiten ist schwierig, da für die untersuchte Aufgabe -- stündliche Prognose der Fahraktivität -- keine etablierten Benchmarks existieren. Die in \cref{sec:verwandte_arbeiten} gesichteten Arbeiten adressieren angrenzende, aber verschiedene Probleme. 


Ein Zweig der Forschungen, welcher eine gewisse Überschneidung mit der vorliegenden Arbeit aufweisen, beschäftigen sich mit dem \textit{Ladezustand} beziehungsweise Ladeverhalten von Elektrofahrzeugen und utilisiert dazu Fahraktivität. 

Der zweite Zweig sagt den nächsten \textit{Aufenthaltsort} voraus, oder gar eine Kette an Orten. Eine der Literaturen, von Lindroth et al., verfolgt den Ansatz, Abfahrtszeiten und Fahrtdistanzen der ersten Reise des Tages vorherzusagen \cite{lindroth2024online}. Die vorliegende Arbeit unterscheidet sich von Literaturen aus diesem Zweig dadurch, dass sie Aktivität selbst als Zielgröße isoliert und den Daten- vom Algorithmuseffekt trennt -- ein methodischer Beitrag, der in den genannten Arbeiten mangels quellenübergreifendem Design nicht möglich ist.

Inhaltlich lassen sich die Ergebnisse dennoch einordnern: Der klare Abstand der Modelle zur Klimatologie-Baseline (\cref{tab:baselines}) bei zugleich quellenabhängiger Obergrenze ist konsistent mit den theoretischen Grenzen der Vorhersagbarkeit menschlicher Mobilität, wie sie in der Literatur von Song et al. beschrieben wird \cite{song2010limits}. Dass die Prognosegüte stärker mit der Datenquelle als mit dem Algorithmus variiert, verschiebt den in der Literatur oft auf die Modellwahl gelegten Fokus hin zur Datenqualität.

% -----------------------------------------------------------------------------
\section{Limitationen}
\label{sec:limitationen}

Durch die Evaluation der Ergebnisse konnten die Forschungsfragen beantwortet werden. Dennoch sind im Laufe der Arbeit Limitationen aufgetreten, die die Ergebnisse beeinflussen können. Die Limitationen betreffen insbesondere die Datenquellen, die Modellwahl und die Evaluation.

Im Laufe der Arbeit sind verschiedene Limiationen aufgetreten, die die Ergebnisse beeinflussen: Heterogene Aktivitätsdefinition, die Datenansprüche des LSTM, Subsampling, heuristische Label-Ableitung und die Eigenheiten des Datensatzes \texttt{VED}.

Die Datenquellen sind in der Arbeit auf fünf Datensätze beschränkt, die unterschiedliche Eigenschaften aufweisen. Ergebnisse können daher nicht immer auf andere Datensätze übertragen werden. Die Modellwahl ist auf drei Verfahren beschränkt, die unterschiedliche Eigenschaften aufweisen. Ergebnisse können daher zwar verglichen werden, aber nicht unbedingt auf andere Verfahren übertragen werden. Die Evaluation ist auf die Metriken F1-Score, ROC-AUC und PR-AUC beschränkt, die unterschiedliche Eigenschaften aufweisen. Ergebnisse können daher zwar verglichen werden, aber nicht unbedingt auf andere Metriken übertragen werden.

Limitationen der Arbeit sind insbesondere die aktive Datenerhebung, was zur folge hat, dass die Datenquellen nicht zwingend repräsentativ für die durchschnittliche deutsche Bevölkerung sind. Realer Verkehr ist ein weiterer Faktor, der einen spürbaren Einfluss auf das individuelle Mobilitätsverhalten hat, da die Datenquellen nicht zwingend repräsentativ für den realen Verkehr sind. 

% -----------------------------------------------------------------------------
\section{Implikationen}
\label{sec:implikationen}

\todo{Welche praktischen und theoretischen Implikationen ergeben sich aus den Ergebnissen?}
